\section{基于KG-GCN的燃煤机组汽轮机系统故障诊断方法}
\label{sec:kg_gcn_v6}

\subsection{引言}

燃煤机组汽轮机系统故障具有多参数耦合和跨设备传播特征。传统基于阈值或经验规则的诊断方法通常依赖固定规则描述异常现象，当机组运行边界变化或故障表现发生偏移时，其适应性受到限制；纯数据驱动分类方法能够从历史样本中学习复杂非线性特征，但模型通常将监测参数视为规则排列的输入变量，难以显式利用设备之间客观存在的热力连接关系。当不同故障共享部分监测参数时，仅依赖局部数据特征容易产生类别混淆。

第三章构建的汽轮机系统故障知识图谱将设备、参数、异常和故障之间的关联组织为图结构，可为多参数故障诊断提供静态领域知识；第二章健康状态模型则能够根据当前运行边界给出设备正常状态参考，并由实测值与健康预测值之间的偏差形成动态状态信息。基于此，本文将知识图谱参数关系与状态监测残差进行融合，采用图卷积网络在参数拓扑上逐层聚合动态特征，并通过无标签残差窗口掩码预训练增强少故障标签条件下的特征学习能力，构建面向汽轮机系统的KG-GCN故障诊断模型。

\subsection{图卷积网络}
\label{subsec:gcn_theory_v6}

复杂工业系统中的设备和参数关系本质上构成非欧几里得拓扑，而传统卷积神经网络主要面向规则网格数据，难以直接表示任意节点之间的结构依赖。图卷积网络（Graph Convolutional Network，GCN）能够依据图的邻接关系聚合节点自身及邻域特征，从而在图结构数据上提取高阶关联信息\cite{Kipf2017GCN}。

设参数图为$G=(V,E)$，其中$V=\{v_1,v_2,\ldots,v_N\}$为节点集合，$E$为节点之间的关联边。图结构由邻接矩阵$A\in\mathbb R^{N\times N}$表示。为在图卷积过程中同时保留节点自身信息，在邻接矩阵中加入单位矩阵$I_N$，得到
\begin{equation}
\tilde A=A+I_N.
\end{equation}
相应度矩阵定义为
\begin{equation}
\tilde D_{ii}=\sum_j\tilde A_{ij}.
\end{equation}
进一步对邻接矩阵进行对称归一化：
\begin{equation}
\bar A=\tilde D^{-\frac{1}{2}}\tilde A\tilde D^{-\frac{1}{2}}.
\end{equation}
设第$l$层节点特征矩阵为$H^{(l)}$，图卷积层的传播形式可写为
\begin{equation}
H^{(l+1)}=\sigma\left(\bar A H^{(l)}W^{(l)}\right),
\end{equation}
式中：$W^{(l)}$为第$l$层可学习权重矩阵；$\sigma(\cdot)$为非线性激活函数。由上式可知，每一层图卷积均按照邻接矩阵规定的连接路径融合目标节点和相邻节点信息，随着层数增加，节点能够逐步获得更远邻域内的结构化特征。

本文在标准GCN基础上采用两层图卷积完成汽轮机参数关联信息聚合。考虑固定DCS节点自身残差同样具有诊断价值，在图卷积层中保留节点自身表示，并引入残差连接和LayerNorm，以减小多层传播过程中节点局部信息的衰减。其计算形式为
\begin{equation}
H^{(1)}=\mathrm{ReLU}\left[\mathrm{LN}\left(H^{(0)}+\bar A H^{(0)}W^{(1)}\right)\right],
\end{equation}
\begin{equation}
H^{(2)}=\mathrm{ReLU}\left[\mathrm{LN}\left(H^{(1)}+\bar A H^{(1)}W^{(2)}\right)\right].
\end{equation}
第一层GCN主要融合直接相邻参数的信息，第二层进一步聚合两跳邻域特征，使节点表示同时包含自身动态变化和局部热力拓扑信息。

\subsection{KG-GCN模型输入构建}
\label{subsec:kg_input_v6}

第三章构建的汽轮机系统故障知识图谱包含设备、参数、异常、故障和处置等多种类型实体，而实时图卷积计算需要将数值状态加载至可观测参数节点。为此，从系统级知识图谱中提取能够与实际DCS量建立映射的参数实体，并依据设备归属、介质流向和直接热力接口关系形成诊断参数图。设任务参数图为
\begin{equation}
G_p=(V_p,E_p),
\end{equation}
式中：$V_p$为可加载运行状态的参数节点集合；$E_p$为由系统知识图谱投影得到的参数关联边。根据$E_p$构建邻接矩阵$A_p$，并按前述方法得到归一化邻接矩阵$\bar A_p$，作为图卷积的静态结构输入。

对于第$i$个状态参数，在时刻$t$的DCS实测值和健康状态模型输出分别记为$y_i(t)$和$\hat y_i(t)$，定义状态偏差
\begin{equation}
e_i(t)=y_i(t)-\hat y_i(t).
\end{equation}
健康状态模型能够解释负荷、入口边界和操作调节造成的大部分正常变化，因此在设备健康状态下，$e_i(t)$主要反映模型预测误差和测量扰动；当设备状态发生异常时，实际运行状态相对于健康参考产生系统性偏离，关联参数残差随之形成持续或协同变化。

由于不同压力和温度参数量纲及幅值范围不同，采用训练集健康残差统计量进行标准化：
\begin{equation}
r_i(t)=\frac{e_i(t)-\mu_i}{\sigma_i+\varepsilon},
\end{equation}
式中：$\mu_i$和$\sigma_i$分别为训练集第$i$个参数健康残差的均值和标准差；$\varepsilon$为防止分母为零的微小常数。验证集和测试集不参与标准化参数估计。

单一时刻的标准化残差可写为
\begin{equation}
\mathbf r_t=[r_1(t),r_2(t),\ldots,r_N(t)]^{\mathrm T}.
\end{equation}
考虑故障演化具有时间连续性，对连续$L_r$个采样时刻构造滑动残差窗口，得到参数--时间残差矩阵
\begin{equation}
R_t=
\begin{bmatrix}
r_1(t-L_r+1)&\cdots&r_1(t)\\
r_2(t-L_r+1)&\cdots&r_2(t)\\
\vdots&\ddots&\vdots\\
r_N(t-L_r+1)&\cdots&r_N(t)
\end{bmatrix}
\in\mathbb R^{N\times L_r}.
\end{equation}
矩阵每一行始终对应知识图谱中的一个固定参数节点，每一列对应连续时间位置。由此，原本独立排列的多参数残差被转换为具有明确节点身份的动态特征矩阵。

在输入层中，首先将每个参数节点的残差时间窗口映射至$d$维初始特征空间：
\begin{equation}
H^{(0)}=\mathrm{ReLU}(R_tW_e+b_e),
\end{equation}
式中：$W_e\in\mathbb R^{L_r\times d}$为特征映射矩阵。随后将$H^{(0)}$与知识图谱得到的归一化邻接矩阵$\bar A_p$共同输入GCN，使参数动态残差沿实际知识关系完成逐层传播。本文固定DCS测点的物理位置具有明确工程意义，因此图级读出阶段保留节点顺序，并将节点局部表示与图传播表示组合后形成图级特征向量$\mathbf g$，用于后续故障分类。

\subsection{KG-GCN模型训练}
\label{subsec:kg_training_v6}

电站长期运行能够持续积累大量未标注运行数据，而高质量故障样本通常数量有限。若仅依赖少量带标签故障样本训练图卷积网络，模型容易过度拟合有限故障模式。为充分利用未标注残差信息，本文采用“残差窗口掩码预训练--固定图卷积编码器--带标签故障分类”的两阶段训练方式，使GCN首先学习参数残差在知识图谱结构中的关联规律，再利用故障标签学习类别判别边界。

设无标签残差窗口集合为
\begin{equation}
\mathcal D_u=\{R^{(1)},R^{(2)},\ldots,R^{(M_u)}\},
\end{equation}
其中$M_u$为无标签窗口数量。预训练阶段对输入残差窗口进行随机遮蔽。设掩码矩阵$M\in\{0,1\}^{N\times L_r}$，$M_{ij}=0$表示第$i$个参数在第$j$个时间位置的残差被遮蔽，$M_{ij}=1$表示该位置保持可见，则遮蔽后的残差窗口为
\begin{equation}
\tilde R=M\odot R+(1-M)\odot R_{mask},
\end{equation}
式中：$\odot$表示逐元素乘法；$R_{mask}$为遮蔽位置的替代值，本文取0。

将遮蔽后的残差窗口与归一化邻接矩阵输入图卷积编码器：
\begin{equation}
Z=E_{\theta}(\bar A_p,\tilde R),
\end{equation}
式中：$E_{\theta}$为待预训练的GCN编码器；$\theta$为其中可学习参数。随后利用解码器将结构化表示恢复至残差窗口空间：
\begin{equation}
\hat R=D_{\phi}(Z),
\end{equation}
式中：$D_{\phi}$为解码器。预训练阶段仅在被遮蔽位置计算重构误差：
\begin{equation}
\mathcal L_{mask}=\frac{\sum_{ij}(1-M_{ij})(R_{ij}-\hat R_{ij})^2}{\sum_{ij}(1-M_{ij})+\varepsilon}.
\end{equation}
通过最小化$\mathcal L_{mask}$，图卷积编码器需要依据可见参数状态及其知识图谱邻域关系恢复被遮蔽残差，从而学习多参数之间的结构关联。

完成掩码预训练后固定图卷积编码器参数，并利用带标签故障窗口训练分类层。设有标签故障样本集合为
\begin{equation}
\mathcal D_l=\{(R^{(n)},y^{(n)})\}_{n=1}^{M_l},
\end{equation}
式中：$R^{(n)}$为第$n$个故障残差窗口；$y^{(n)}$为对应故障类别标签；$M_l$为有标签故障样本数量。故障窗口经过预训练编码器后得到图级特征
\begin{equation}
\mathbf g^{(n)}=E_{\theta^*}(\bar A_p,R^{(n)}),
\end{equation}
其中$\theta^*$表示预训练后固定的图卷积参数。分类层采用全连接映射和Softmax函数输出类别概率：
\begin{equation}
\hat{\mathbf p}^{(n)}=\mathrm{Softmax}\left(W_c\mathbf g^{(n)}+b_c\right).
\end{equation}
对于$C$类故障，分类阶段采用交叉熵损失函数
\begin{equation}
\mathcal L_{cls}=-\frac{1}{M_l}\sum_{n=1}^{M_l}\sum_{k=1}^{C}y_{nk}\log \hat p_{nk}.
\end{equation}
通过上述训练过程，模型利用无标签残差窗口学习知识拓扑约束下的状态关联，再由少量带标签故障样本完成类别边界学习。完成训练后，新的运行残差窗口与参数邻接矩阵共同输入KG-GCN，即可输出对应的故障类别概率。

\subsection{本章小结}

本章构建了基于知识图谱和图卷积网络的燃煤机组汽轮机系统故障诊断方法。首先介绍图卷积网络的节点特征传播过程，并在标准GCN基础上采用两层图卷积和残差连接提取参数邻域信息；随后从第三章系统级知识图谱中抽取可与DCS状态量建立映射的参数关系，将第二章健康状态模型产生的标准化残差按照固定参数节点构造成滑动窗口矩阵，与知识邻接矩阵共同形成KG-GCN输入；在模型训练阶段，通过随机残差遮蔽进行无标签预训练，并在预训练完成后固定图卷积编码器，利用带标签故障样本训练分类层。由此形成“知识图谱参数关系--状态监测残差--图卷积特征传播--掩码预训练--故障分类”的诊断流程，第五章将选取VHP作为代表性设备对所提方法进行算例验证。